\documentclass[aspectratio=169]{beamer}

%\usepackage{graphics}
\usepackage{svg}
\usepackage{booktabs}
\usepackage{tabularx}

% modern minimalist latex theme
\usetheme{metropolis}

% fill blocks with background color
\metroset{block=fill}

% footer
\setbeamertemplate{frame footer}{\insertshortauthor~(\insertshortinstitute)\hfill\insertshorttitle\hfill}
%\setbeamertemplate{frame footer}{\insertshortauthor~(\insertshortinstitute)}


\title[Fidelity of Explanations for AD Classification from MRI]{Evaluating the Fidelity of Explanations for Convolutional Neural Networks in Alzheimer’s Disease Detection}

\author[Hiller]{Bjarne C. Hiller}
\date{2025-03-09}
\institute[Uni Rostock]{University of Rostock}

%\logo{\includegraphics[height=0.5cm]{logos/logo-uni-hro}}

\titlegraphic{
    \includegraphics[height=1cm]{logos/logo-uni-hro}\hfill
    \includegraphics[height=1cm]{logos/logo-vac}
	%\includegraphics{logos/logo-dzne}
}

\begin{document}

\maketitle

% \begin{frame}{Agenda}
% 	\tableofcontents
% \end{frame}

%
\begin{frame}{Who are we?}
	% introduce Martin Becker, Martin Dyrba, and Me
	\begin{block}{University of Rostock: Institute for Visual and Analytic Computing}
		\begin{tabularx}{\textwidth}{XXX}
			Bjarne C. Hiller & Martin Becker & Thomas Kirste \\
		\end{tabularx}
		% 		\begin{itemize}
		% 			\item Bjarne C. Hiller, MSc
		% 			\item Jun.-Prof. Dr. Martin Becker
		% 			\item Prof. Dr.-Ing. Thomas Kirste
		% 		\end{itemize}
	\end{block}

	\begin{block}{German Center for Neurodegenerative Diseases (DZNE)}
		\begin{itemize}
			\item Devesh Singh, MSc
			\item Dr. Martin Dyrba
		\end{itemize}
	\end{block}
\end{frame}

\begin{frame}{Deep Learning for Medical Image processing}
	% introduce deep learning and XAI
	% - want to analyse MRI from AD subjects
	% - Deep Learning powerful for image processing, expert level performance
	% - however, Black Box systems
	% - Can we trust prediction?
\end{frame}

\begin{frame}{Data}
	% - ADNI data set overview
	% - preprocessing
\end{frame}

\begin{frame}{Model}

\end{frame}

\begin{frame}{eXplainable AI and Feature Attribution Methods}
	% - explain overall principle of feature attribution methods
	% - plethora of feature attribution methods
	% - Can we trust *explanation*?
	% - "Not sure, if model is broken, or explanation is wrong"
\end{frame}

\begin{frame}
	\begin{center}
		\includegraphics[width=0.5\textwidth]{figures/fry-and-xai.jpg}
	\end{center}
\end{frame}

\begin{frame}{Perturbation Tests: Insertion and Deletion}
	% introduce perturbation tests
\end{frame}

\begin{frame}{Attribution Maps}

\end{frame}

\begin{frame}{Relevance per ROI}

\end{frame}

\begin{frame}{results}
	% results perturbation test
\end{frame}

\begin{frame}{Take-Away}
	\begin{enumerate}
		\item Perturbation tests offer a model-agnostic fidelity metric.
		\item The baseline should be chosen carefully w.r.t. to context.
		\item Attribution Maps need interpretation to actually explain anything.
	\end{enumerate}
\end{frame}

\end{document}